\scanpage{5}
si
\[
p_*=(p_1,\ldots,p_n)\in([3,+\infty]\cap\overline{\mathbb{N}})^n
\]
\deleted{$\unsure{\underline{R}=\Gamma}$}
\[
\Gamma=\Gamma_{p_*}
\]
est entièrement déterminé par $p_*$ et, quel que soit $p_*$, les conditions \deleted{$\unsure{\sigma_J\neq1|\sigma_I\neq1}\text{ dans }\Gamma_{p_*}$}
\[
\sigma_J\neq1\quad\text{dans }\Gamma_{p_*}
\]
sont satisfaites dans $\Gamma=\Gamma_{p_*}$. Tous les $\underline{R}$ sont isom. entre eux, et quand on prend des $\underline{R}$ munis d'un ``repère'' $\alpha\in R$, ils sont canoniquement isom. entre eux, et \inserted{can.} isom. à \deleted{$\unsure{E|R}$}
\[
\underline{R}_{p_*}=G_n/\Gamma_{p_*}
\]
munis du repère image de $1\in G_n$.

Pour tte partie $H$ de \deleted{$\unsure{G}$} $[0,n]\cap\mathbb{N}$, on désigne par $G_{nH}$ le sous-groupe de $G_n$ engendré par les $\sigma_i$, $i\notin H$. Donc
\[
H\subset K\longrightarrow G_K\subset G_H
\]
\diagram{\begin{tikzcd}[ampersand replacement=\&]
\underline{R}/G_K \arrow[r] \arrow[d,equal] \& \underline{R}/G_H \arrow[d,equal]\\
D_K \& D_H
\end{tikzcd}}
Si $H=\{i\}$, on trouve
\[
\begin{aligned}
D_{\{i\}}=D_i&=\text{ens des }i\text{-facettes}\\
&\text{du }\unsure{p|n}\text{-polyèdre régulier }\Pi\\
&\text{défini par le}
\end{aligned}
\]
% archive sheet number: 4
